\section{An unbounded family with prescribed local level data}
\label{sec:powerfree-global-family}

The preceding local constructions can be supplied by rational elliptic
curves of unbounded height while the relevant local bounding domain remains
fixed. We prove this using a least-prime theorem and a finite sieve in one
arithmetic progression. The result is an existence family with specified
reduction, torsion-image, and level-window properties. No lower bound for
its abc quality with an exponent greater than one is asserted.

\subsection{A uniform finite sieve}

An integer is \(k\)-power-free if it is not divisible by \(q^k\) for any
prime \(q\). We shall use the following estimate with \(k\) and the
progression modulus both growing.

\begin{lemma}[Power-free values in a progression]
\label{lem:powerfree-crt-count}
Let \(k\ge5\), let \(M\) be a positive even integer, and let \(X\ge1\).
Fix a residue class \(r\bmod M\), and suppose that, for every integer
\(A\in[X,2X]\) in this class, none of
\[
 A,\qquad A^2+1,\qquad 2A^2+1
\]
is divisible by \(q^k\) at a prime \(q\mid M\). The number \(N_{\rm good}\)
of such \(A\) for which all three values are \(k\)-power-free satisfies
\begin{equation}\label{eq:pf-sieve}
 N_{\rm good}\ge
 \frac XM\left(1-5\sum_{q\ {\rm prime}}q^{-k}\right)
       -1-5(9X^2)^{1/k}
 \ge \frac{3X}{4M}-1-5(9X^2)^{1/k}.
\end{equation}
\end{lemma}

\begin{proof}
The interval contains at least \(X/M-1\) integers of the prescribed
class. If \(q\nmid M\), then \(q\) is odd. The polynomial \(T\) has one
root modulo \(q^k\), and each of \(T^2+1\) and \(2T^2+1\) has at most two.
For the latter assertion, there are at most two roots modulo \(q\), and
the derivative is nonzero at every root. A root \(x\bmod q^h\) has lifts
\(x+tq^h\); reduction of
\[
 f(x+tq^h)\equiv f(x)+tq^h f'(x)\pmod {q^{h+1}}
\]
gives exactly one admissible \(t\bmod q\). This proves the root bound
at every exponent without a uniformity assumption on the prime.

Each of these at most five roots, combined with \(r\bmod M\), specifies
one class modulo \(Mq^k\). It contributes at most \(X/(Mq^k)+1\)
integers. All three polynomial values are positive and at most \(9X^2\).
Thus only primes
\[
 q\le Y:=(9X^2)^{1/k}
\]
can exclude an integer. Summing these upper bounds over \(q\le Y\)
excludes at most
\[
 \frac{5X}{M}\sum_{q\ {\rm prime}}q^{-k}+5Y
\]
integers. The sum of the error terms is finite; no error of size one
is summed over infinitely many primes. Subtracting gives the first
inequality. Finally,
\[
 \sum_{q\ {\rm prime}}q^{-k}
 \le \sum_{n=2}^{\infty}n^{-k}
 \le 2^{-k}+\frac{2^{1-k}}{k-1}
 \le \frac32\,2^{-k}.
\]
For \(k\ge5\), five times this bound is at most \(15/64<1/4\).
\end{proof}

\subsection{The existence theorem}

The only prime-distribution input needed below is the effective estimate
of Xylouris \cite[Theorem~1.1, author manuscript p.~9]{XylourisLinnik2009}.
With \(P(m)\) the maximum, over reduced classes modulo \(m\), of the least
prime in the class, it gives an effective absolute constant \(C\) such
that
\begin{equation}\label{eq:pf-linnik}
 P(m)\le C m^{26/5}.
\end{equation}
We use the exponent \(5.2\) of that archived author version throughout;
the later improvement is unnecessary. In particular, \(C\) is independent
of both the modulus and the reduced class.

\begin{theorem}[An unbounded balanced Frey--Legendre family]
\label{thm:powerfree-global-family}
For every sufficiently large prime \(\ell\equiv43\pmod {60}\), let \(p\)
be the least prime congruent to \(-1\pmod {30\ell}\), and put
\begin{equation}\label{eq:pf-parameters}
 M=10\ell p^2,\qquad X=\exp(\ell^2/16).
\end{equation}
There are at least \(X/(2M)\) integers \(A\in[X,2X]\) satisfying
\begin{equation}\label{eq:pf-crt}
 A\equiv1\pmod2,\qquad A\equiv1\pmod {5\ell},
 \qquad A\equiv p\pmod {p^2},
\end{equation}
for which \(A,A^2+1,2A^2+1\) are all \(\ell\)-power-free.
Every such \(A\) gives the positive primitive triple and curve
\begin{equation}\label{eq:pf-family}
 (a,b,c)=(A^2,A^2+1,2A^2+1),\qquad
 D_A:\ y^2=x(x-a)(x+b).
\end{equation}
The curve is rationally isomorphic to the Legendre curve with
\[
 \lambda_A=-1-A^{-2}.
\]
For \(L_A=\Q(i,D_A[30])\), the following assertions hold:
\begin{enumerate}
\item At \(p\), the curve is split multiplicative with native Tate
      valuation \(v_p(q)=4\).
\item The mod-\(\ell\) image over \(\Q\) and over \(L_A\) contains
      \(\operatorname{SL}_2(\mathbf F_\ell)\). Reduction at \(\ell\) is good.
\item Every nonzero integral Tate order over \(L_A\) is prime to \(\ell\).
\item The Tate degree normalized by \([L_A:\Q]^{-1}\) is
      \begin{equation}\label{eq:pf-degree}
       Q_A=2\log\frac{A^2(A^2+1)(2A^2+1)}2,
       \qquad \frac34\ell^2<Q_A<\ell^2.
      \end{equation}
\end{enumerate}
The parameters lie in one fixed compactly bounded domain in the original
sense specified below, and both \(\height(\lambda_A)\) and
\(\height(j(D_A))\) tend to infinity with \(\ell\). The threshold is
effective in terms of the effective constant in \eqref{eq:pf-linnik};
no numerical value of that constant is extracted here.
\end{theorem}

\begin{proof}[Existence of the integers]
The three moduli in \eqref{eq:pf-crt} are pairwise coprime. By
\eqref{eq:pf-linnik},
\[
 30\ell-1\le p\le C(30\ell)^{26/5},\qquad
 M\le C_0\ell^{57/5},\qquad C_0=10C^2 30^{52/5}.
\]
The primes dividing \(M\) cause no exclusion in
Lemma~\ref{lem:powerfree-crt-count}. At 2, \(A\) is odd, so
\(v_2(A^2+1)=1\) and \(2A^2+1\) is odd. At 5 and \(\ell\), the three
values are \(1,2,3\). At \(p\), their valuations are \(1,0,0\).
All four primes are distinct.

For \(k=\ell\), the error parameter is
\[
 Y=9^{1/\ell}\exp(\ell/8).
\]
It suffices that \(1+5Y\le X/(4M)\). Since
\(1+5Y\le46\exp(\ell/8)\), an explicit sufficient condition is
\begin{equation}\label{eq:pf-threshold}
 \frac{\ell^2}{16}-\frac{\ell}{8}
       -\frac{57}{5}\log\ell\ge\log(184C_0).
\end{equation}
The left side tends to infinity and is increasing for \(\ell\ge43\):
its derivative is \(\ell/8-1/8-57/(5\ell)>0\).
Lemma~\ref{lem:powerfree-crt-count} now gives the lower bound \(X/(2M)\).
All constants are independent of the chosen \(\ell,p,A\).

There are infinitely many eligible primes \(\ell\). Indeed, if the primes
congruent to 43 modulo 60 were a finite list \(\ell_1,\ldots,\ell_s\),
the reduced CRT class congruent to 43 modulo 60 and to 1 modulo each
\(\ell_i\) would contain a prime by \eqref{eq:pf-linnik}. It would be a
new prime in the original progression, a contradiction.
\end{proof}

\begin{proof}[Arithmetic and height assertions]
Pairwise coprimality in \eqref{eq:pf-family} follows from
\(\gcd(A^2,A^2+1)=1\) and \(2(A^2+1)-(2A^2+1)=1\).
The substitution \(x=A^2x_1,\ y=A^3y_1\) gives
\[
 y_1^2=x_1(x_1-1)(x_1+1+A^{-2}).
\]
It is an isomorphism over \(\Q\), not a quadratic twist.
The general invariants \eqref{eq:fr-invariants} give, with
\(S_A=3A^4+3A^2+1\),
\begin{equation}\label{eq:pf-invariants}
 \Delta=16(abc)^2,\qquad c_4=16S_A,\qquad
 j(D_A)=\frac{256S_A^3}{(abc)^2},\qquad \gcd(S_A,abc)=1.
\end{equation}
At an odd prime \(r\mid abc\), unit \(c_4\) makes the equation minimal
and multiplicative, with Tate order \(2v_r(abc)\).
Splitness is also uniform in \(A\). If \(r\mid a\) or \(r\mid b\),
the reduced node has tangent slopes \(\pm1\); if \(r\mid c\), translation
of the node at \(x=a\) gives \(y^2=x_1^2(x_1+a)\), with slopes \(\pm A\).
In each case they are distinct units.
These are the reduction criteria of
\cite[Chapter~VII, Proposition~5.1]{Silverman}; the split Tate
description is \cite[Theorems~1--2]{KedlayaTate}.
At the prescribed \(p\), \(v_p(A)=1\), and hence \(v_p(q)=4\).
At \(\ell\), the endpoints are \(1,2,3\), so reduction is good.

The good reduction at 5 is the same fixed curve
\(y^2=x(x-1)(x+2)\) used in Section~\ref{sec:frey-realizations}.
It has four \(\mathbf F_5\)-points, giving Frobenius polynomial
\(T^2-2T+5\). Its discriminant \(-16\) is a nonsquare modulo \(\ell\),
as \(\ell\equiv3\pmod4\); thus this image element fixes no line.
At \(p\), Tate uniformization supplies a nonidentity inertia
transvection of order \(\ell\), because \(p\ne\ell\) and
\(\ell\nmid v_p(q)=4\)
\cite[Proposition~3]{KedlayaTate}.
Conjugating its fixed line by the Frobenius element gives a distinct
fixed line. Lemma~\ref{lem:fr-sl2} therefore gives the full
\(\operatorname{SL}_2(\mathbf F_\ell)\) subgroup.
The same lemma applies over \(L_A\): the latter is Galois, and
\eqref{eq:fr-level-degree} gives
\begin{equation}\label{eq:pf-level-degree}
 [L_A:\Q]\mid 2\lvert\GL_2(\Z/30\Z)\rvert
       =276480=2^{11}3^3 5,
\end{equation}
which is prime to \(\ell\). The normal image over \(L_A\) retains the
order-\(\ell\) transvection and its conjugate.

At 2, \eqref{eq:pf-invariants} gives \(v_2(j)=6\), hence potential good
reduction \cite[Chapter~VII, Proposition~5.5]{Silverman}. We check why
the specific field \(L_A\) already gives good reduction. Potential good
reduction makes the inertia image on \(T_3(D_A)\) finite. Since all
3-torsion is rational over \(L_A\), that image lies in
\(1+3\Mat_2(\Z_3)\), a torsion-free group. To see the last assertion,
write a nonidentity element as \(1+3^hB\), where \(h\ge1\) and
\(B\not\equiv0\pmod3\). Powers prime to 3 preserve \(h\), and
\[
 (1+3^hB)^3
 =1+3^{h+1}\bigl(B+3^hB^2+3^{2h-1}B^3\bigr)
\]
increases it by exactly one. Thus no positive finite power is the
identity. Inertia is trivial, and the Neron--Ogg--Shafarevich criterion
\cite[Chapter~VII, Theorem~7.1]{Silverman} gives good reduction over
every completion of \(L_A\) at 2.

For \(w\mid r\), where \(r\) is odd and divides \(abc\), the integral
Tate order over \(L_A\) is
\[
 \operatorname{ord}_w(q_w)=2e(w/r)v_r(abc).
\]
If \(r\mid A\), this equals \(4e(w/r)v_r(A)\); otherwise it is twice
the ramification index times the valuation of one of the two quadratic
factors. Each latter valuation, and \(v_r(A)\), lies strictly between
zero and \(\ell\). Also \(\ell\nmid e(w/r)\), by
\eqref{eq:pf-level-degree}. Since \(\ell\) is odd, every nonzero order
is prime to \(\ell\). We do not need or assert that \(A^2\) itself is
\(\ell\)-power-free.

Sum these orders with residue-degree weights and divide by \([L_A:\Q]\).
For each \(r\), the identity \(\sum_{w\mid r}e(w/r)f(w/r)=[L_A:\Q]\),
good reduction at 2, and \(v_2(abc)=1\) give the exact formula for \(Q_A\)
in \eqref{eq:pf-degree}. Finally, for \(A>1\),
\[
 A^6<\frac{A^2(A^2+1)(2A^2+1)}2<3A^6.
\]
Together with \(X\le A\le2X\), this yields
\begin{equation}\label{eq:pf-height-window}
 \frac34\ell^2<Q_A
 <\frac34\ell^2+12\log2+2\log3
 <\frac34\ell^2+16<\ell^2.
\end{equation}
The remaining domain and unbounded-height assertions follow from
Proposition~\ref{prop:powerfree-fixed-domain} below.
\end{proof}

The prime-to-\(\ell\) order assertion is over \(L_A\), as stated. It is
not an assertion over \(L_A(D_A[\ell])\), whose ramification at \(p\)
has an additional factor \(\ell\).
Set \(\delta=2^{12}3^3 5=552960\), the constant for \(d_{\rm mod}=1\)
in \cite[Theorem~5.7.1]{JoshiIV}. Equation~\eqref{eq:pf-height-window}
also proves both numerical endpoints
\begin{equation}\label{eq:pf-joshi-window}
 \sqrt{Q_A}<\ell
 <10\delta\sqrt{Q_A}\log(2\delta\log Q_A).
\end{equation}
For the second inequality, use \(\sqrt{Q_A}>\ell/2\) and
\(\log(2\delta\log Q_A)>1\). This directly verifies a window for our
chosen \(\ell\); it does not identify it with an existential choice in
that source.

\subsection{The full local torsion field}

The local field in this family is determined with its actual Tate unit
retained. Put \(N=30\ell\) and \(e=15\ell\). Full rational 2-torsion
implies that a square root \(b_0\) of \(q\) belongs to \(\Q_p\).
Indeed, its Tate class is rational, so for every Galois element the ratio
of a conjugate to \(b_0\) is a power of \(q\). The ratio has valuation
zero, hence is one. Thus \(v_p(b_0)=2\).

Let \(K_0=\Q_p(\mu_N)\), choose \(\varpi^e=b_0\), and write
\(u_0=b_0/p^2\). Since \(p\equiv-1\pmod N\), \(K_0/\Q_p\) is
unramified quadratic. The element \(\varpi\) has valuation \(2/e\);
the uniformizer is instead
\begin{equation}\label{eq:pf-uniformizer}
 \beta=\frac{\varpi^{(e+1)/2}}p,\qquad
 \beta^e=p\,u_0^{(e+1)/2},\qquad
 \varpi=u_0^{-1}\beta^2.
\end{equation}
The equation for \(\beta\) is Eisenstein over \(K_0\). Consequently
\[
 E=K_0(\varpi)=K_0(\beta),\qquad
 [E:\Q_p]=30\ell,\quad e(E/\Q_p)=15\ell,\quad f(E/\Q_p)=2.
\]
Tate uniformization identifies
\[
 \Q_p(D_A[N])=\Q_p(\mu_N,q^{1/N})=E:
\]
fixing each torsion class forces its root-of-unity multiplier to be one,
because a nontrivial power of \(q\) has nonzero valuation.
The unramified quadratic field \(K_0\) contains \(i\), so \(E\) is also
the completion at \(p\) of \(\Q(i,D_A[N])\).
It is Galois, being the splitting field obtained from the cyclotomic
field by adjoining all roots of \(T^e-b_0\).
Moreover,
\begin{equation}\label{eq:pf-small-tame}
 e=15\ell\le p-2,\qquad \gcd(e,p-1)=1,
\end{equation}
since \(p\ge30\ell-1\) and \(p-1\equiv-2\pmod {15\ell}\).
Thus the small tame range of the local arguments occurs at unbounded
global height. No substitution \(q=p^4\) or \(b_0=p^2\) has been made.

\subsection{A fixed domain and the finite-set quantifier}

Let \(U=\mathbf P^1\setminus\{0,1,\infty\}\). Consider
\begin{equation}\label{eq:pf-domains}
 \begin{split}
  \mathcal K_\infty
     &=\{z\in\mathbf C:|z+1|\le1/2\},\\
  \mathcal K_2
     &=\{z\in\overline{\Q}_2:|z+2|_2\le1/8\}.
 \end{split}
\end{equation}
For nonarchimedean places, the original compactly bounded condition of
\cite[Example~1.3(ii), author manuscript pp.~5--6]{MochGeneralPosition}
is formulated on intersection with each finite local extension. It does
not require the entire set \(\mathcal K_2\) to be compact in
\(\overline{\Q}_2\).

\begin{proposition}[Fixed local bounds and unbounded rational height]
\label{prop:powerfree-fixed-domain}
The sets \eqref{eq:pf-domains} give one compactly bounded domain of support
\(\{\infty,2\}\) in that original sense. Every member of
Theorem~\ref{thm:powerfree-global-family} belongs to this domain, and
\begin{equation}\label{eq:pf-unbounded-heights}
 \height(\lambda_A)=\log(A^2+1)\longrightarrow\infty,\qquad
 \height(j(D_A))\ge Q_A\longrightarrow\infty.
\end{equation}
On all of \(\mathcal K_2\), the Legendre invariant has valuation 6.
\end{proposition}

\begin{proof}
Both sets avoid the three removed points and are stable under the
respective Galois actions. The complex set is a compact regular closed
disc. For every finite extension \(k/\Q_2\),
\[
 \mathcal K_2\cap k=-2+8\OO_k
\]
is compact and open, hence is the closure of its interior in \(k\).
The set \(\mathcal K_2\) is itself open and closed in
\(\overline{\Q}_2\), so any separate regular-closed requirement also
holds. These are the stated original conditions.

For \(A\ge2\), \(|\lambda_A+1|=A^{-2}\le1/4\). Since \(A\) is odd,
\(A^2\) and its inverse are 1 modulo 8, giving
\(\lambda_A\equiv-2\pmod8\). Thus the domain is fixed independently of
\(\ell,p,A\), and the degree bound is one. If \(z\in\mathcal K_2\),
then \(v_2(z)=1\), \(v_2(1-z)=0\), and
\(v_2(1-z+z^2)=0\). The formula
\[
 j(z)=\frac{256(1-z+z^2)^3}{z^2(1-z)^2}
\]
therefore gives \(v_2(j(z))=6\).
On \(\mathcal K_\infty\), the inequalities
\(1/2\le|z|\le3/2\) and \(|1-z|\ge3/2\) also give the fixed coarse bound
\[
 |j(z)|\le
 \frac{256(19/4)^3}{(1/2)^2(3/2)^2}
 =\frac{438976}{9}<50000.
\]

The numerator and denominator of \(\lambda_A=-(A^2+1)/A^2\)
are coprime. Its rational height is therefore the first expression in
\eqref{eq:pf-unbounded-heights}. Also \(abc/2\) is odd, and
\eqref{eq:pf-invariants} reduces to
\[
 j(D_A)=\frac{64S_A^3}{(abc/2)^2}
\]
in lowest terms. Its denominator alone gives
\(\height(j(D_A))\ge2\log(abc/2)=Q_A\).
\end{proof}

\begin{corollary}[Avoidance of any fixed finite set]
\label{cor:powerfree-finite-exception}
If a finite set of parameters in \(U(\overline{\Q})\), or a finite set
of geometric elliptic moduli points, is fixed before \(\ell\) varies,
all sufficiently large members of the family avoid it.
\end{corollary}

\begin{proof}
Take the maximum of the finitely many parameter heights, or of the
heights of their \(j\)-invariants, and apply
\eqref{eq:pf-unbounded-heights}. In particular, bounded complex absolute
values of these invariants do not bound their rational heights.
\end{proof}

\subsection{The source-interface limits}

The finite-set statement has its exact order of quantifiers: the set
must be fixed before the family varies. It cannot be used for an
exceptional set chosen separately for each member. There is also a
topological difference that prevents silently substituting one definition
for another.

\begin{proposition}[The literal ambient compactness condition]
\label{prop:powerfree-literal-compact}
In the usual valuation topology of \(U(\overline{\Q}_2)\), every compact
subset has empty interior. Consequently no nonempty compact subset is
the closure of its ambient interior.
\end{proposition}

\begin{proof}
Every nonempty open ball contains, after translation and a sufficiently
small nonzero scaling, infinitely many algebraic integral units with
distinct residues in \(\overline{\mathbf F}_2\). For example, one may
take lifts in finite unramified extensions. The distances between the
scaled points are all the same positive constant. Thus the ball is not
totally bounded and cannot be contained in a compact set. A compact
set with nonempty interior would contain such a ball, a contradiction.
\end{proof}

Section~5.1 of the archived v2 of \cite{JoshiIV}, p.~50, literally
requires nonempty compact subsets in the product of the local algebraic
closures and equality with the closure of their ambient interiors.
Under the usual valuation topology, the proposition applies to the
nonarchimedean factor, and the same obstruction holds for a product
containing it. The domain \eqref{eq:pf-domains} satisfies the original
finite-extension-section definition; it must not be described as a
compact subset of that full algebraic closure. Using finite-extension
sections or a different topology would require stating that interpretation
explicitly. This observation does not disprove the original IUT statements.

Finally, the separate existence quantifier in
\cite[Theorem~5.7.1, p.~53]{JoshiIV} supplies some prime in a window.
Our chosen \(\ell\) satisfies \eqref{eq:pf-joshi-window} and the arithmetic
properties proved above directly, but these facts do not identify it
with that existential prime. The covering and marking requirements of
\cite[Definition~3.1]{IUT1} for these members are supplied separately
by Corollary~\ref{cor:itd-powerfree-family}. They do not identify the
complete possible-image source including subsequent indeterminacies.
The classical least-prime,
elliptic-curve, and local-field inputs in this section remain mathematical
source dependencies; they are not asserted as new Lean axioms or as a
formal proof of the abc conjecture.
